\begin{abstract}
\noindent The radius-period diagram of the NASA Exoplanet Archive is widely read as a map of where planets are, yet its radii come from two routes: transit depth, and a mass-radius relation applied to planets without a transit. We ask how catalogued planets are distributed in radius and orbital period, and how that distribution changes when the two routes are separated using the archive's own label. We describe 6,045 planets with both values.
% src: estimation_results.json#sample.n_observations
Planets with radii calculated from mass, about a fifth of the sample, have a median radius about five times and a median period about 45 times those of transit-measured planets.
% src: estimation_results.json#comparisons.calculated_minus_measured_transiting.lp_median_diff.ratio_10_pow
The small-planet hump consists mostly of transit-measured radii, while more than half of the giant-planet hump consists of calculated radii, which the assumed relation caps at 14.3 Earth radii.
% src: estimation_results.json#statistics.radius_calculated.max
The two provenances therefore describe different parts of the catalogue, and pooled patterns in the diagram mix them.

\medskip
\noindent Keywords: exoplanets, planet radius, orbital period, NASA Exoplanet Archive, mass-radius relation.
\end{abstract}
